% from ega0-20-fr.tex
\oldpage[0\textsubscript{IV-1}]{146}
que $\mathfrak{R}'$ soit un idéal de $B'$ tel que
$f(\mathfrak{R})\subset\mathfrak{R}'$, et que
$E'=B'/\mathfrak{R}'^2$ soit une $B'$-extension $\Lambda'$-triviale de
$C'=B'/\mathfrak{R}'$ par $\mathfrak{R}'/\mathfrak{R}'^2$. On a alors un
diagramme commutatif
\begin{equation}
\tag{20.6.22.2}
\label{0.20.6.22.2-fr}
\xymatrix@C=0.7em{
  0 \ar[r] & \Upsilon^{C'}_{B'/A'/\Lambda'} \ar[r] &
  \Upsilon_{C'/A'/\Lambda'} \ar[r] &
  \mathfrak{R}'/\mathfrak{R}'^2 \ar[r] &
  \Omega^1_{B'/A'}\otimes_{B'} C' \ar[r] &
  \Omega^1_{C'/A'} \ar[r] & 0 \\
  0 \ar[r] & \Upsilon^C_{B/A/\Lambda} \ar[r] \ar[u] &
  \Upsilon_{C/A/\Lambda} \ar[r] \ar[u] &
  \mathfrak{R}/\mathfrak{R}^2 \ar[r] \ar[u] &
  \Omega^1_{B/A}\otimes_B C \ar[r] \ar[u] &
  \Omega^1_{C/A} \ar[r] \ar[u] & 0
}
\end{equation}
comme il résulte de \sref{0.20.6.21.3-fr} et \sref{0.20.5.11.3-fr}.

\begin{corollary}[20.6.23]
\label{0.20.6.23-fr}
Sous les hypothèses de \sref{0.20.6.22-fr}, supposons en outre que $B$ soit
une $A$-algèbre formellement lisse. Alors on a une suite exacte
\begin{equation}
\tag{20.6.23.1}
\label{0.20.6.23.1-fr}
  0\to\Upsilon_{C/A/\Lambda}\xrightarrow{\chi_E}
  \mathfrak{R}/\mathfrak{R}^2\xrightarrow{\delta_{C/B/A}}
  \Omega^1_{B/A}\otimes_B C\to\Omega^1_{C/A}\to0.
\end{equation}
\end{corollary}

\begin{proof}
Cela résulte en effet de \sref{0.20.6.18-fr}, (ii).
\end{proof}

\begin{env}[20.6.24]
\label{0.20.6.24-fr}
Lorsque les hypothèses de \sref{0.20.6.22-fr} sont satisfaites, on dit encore
que l'homomorphisme caractéristique $\chi_E$ est
l'\emph{homomorphisme caractéristique de la $A$-algèbre $B$, relativement à
$\Lambda$ et à l'idéal $\mathfrak{R}$} (en supprimant ces dernières
précisions lorsqu'il n'y a pas de confusion à craindre) ; on le notera
parfois $\chi_B$ ou $\chi_{B/A}$.
\end{env}

\begin{proposition}[20.6.25]
\label{0.20.6.25-fr}
Soient $s:\Lambda\to A$, $u:A\to B$, $v:B\to C$ trois homomorphismes
d'anneaux, $L$ un $C$-module. On a alors une suite exacte
\begin{equation}
\tag{20.6.25.1}
\label{0.20.6.25.1-fr}
\begin{split}
  0\to\operatorname{D\acute{e}r}_B(C,L)\xrightarrow{u^1}
  \operatorname{D\acute{e}r}_A(C,L)\xrightarrow{v^1}
  \operatorname{D\acute{e}r}_A(B,L)\xrightarrow{\partial} \\
  \to\operatorname{Exalcom}_{B/\Lambda}(C,L)\xrightarrow{u^1}
  \operatorname{Exalcom}_{A/\Lambda}(C,L)\xrightarrow{v^1}
  \operatorname{Exalcom}_{A/\Lambda}(B,L)\to0
\end{split}
\end{equation}
où $u^1$, $v^1$ sont les homomorphismes définis dans (18.3.6.4) et
(18.3.6.2), et $\partial$ est défini comme dans \sref{0.20.2.2-fr}.
\end{proposition}

\begin{proof}
En effet, comme la suite exacte \sref{0.20.6.16.1-fr} est scindée, on en
déduit une suite exacte
\[
  0\to\operatorname{Hom}_C(\mathrm{K}_\bullet(C/B/\Lambda),L)\to
  \operatorname{Hom}_C(\mathrm{K}'_\bullet(C/A/\Lambda),L)\to
  \operatorname{Hom}_C(\mathrm{K}^C_\bullet(B/A/\Lambda),L)\to0.
\]
Si on applique à ce complexe la suite exacte de cohomologie, en tenant compte
de \sref{0.20.6.15.4-fr}, et de \sref{0.20.6.6-fr}, on obtient
\sref{0.20.6.25.1-fr} car on a
\[
  \operatorname{Hom}_C(\mathrm{K}^C_\bullet(B/A/\Lambda),L)
  =\operatorname{Hom}_B(\mathrm{K}_\bullet(B/A/\Lambda),L)
\]
par définition ; l'identification de $u^1$ et $v^1$ avec les homomorphismes
de (18.3.4.2) résulte de \sref{0.20.6.6.4-fr}.
\end{proof}

\begin{remark}[20.6.26]
\label{0.20.6.26-fr}
Dans ce numéro, les complexes $\mathrm{K}_\bullet(C/B/A)$ sont apparus comme
un artifice technique destiné à simplifier l'exposition de certains
comportements fonctoriels. En réalité, ces complexes, considérés comme objets
de la catégorie des
